\section{模型的建立}

\subsection{问题一分析}
题目要求建立模型描述折叠桌的动态变化图，由于在折叠时用力大小的不同，我们不能描述在某一时刻折叠桌的具体形态，但我们可以用每根木条的角度变化来描述折叠桌的动态变化。首先，我们知道折叠桌前后左右对称，我们可以运用几何知识求出四分之一木条的角度变化。最后，根据初始时刻和最终形态两种状态求出桌腿木条开槽的长度


\subsection{算法示例}

数学建模求解算法示例：
\begin{center}
\begin{minipage}{0.8\textwidth}
\begin{algorithm}[H]%[!htp]
\caption{算法的名字} %算法的名字
{\bf 输入:} %算法的输入， \hspace*{0.02in}用来控制位置，同时利用 \\ 进行换行
input parameters A, B, C\\
{\bf 输出:} %算法的结果输出
output result
\begin{algorithmic}[1]
\State some description 算法介绍 % \State 后写一般语句
\For{condition} % For 语句，需要和EndFor对应
  \State ...
  \If{condition} % If 语句，需要和EndIf对应
    \State ...
  \Else
    \State ...
  \EndIf
\EndFor
\While{condition} % While语句，需要和EndWhile对应
  \State ...
\EndWhile
\State \Return result
\end{algorithmic}
\end{algorithm}
\end{minipage}
\end{center}
\vspace{2ex}
